\documentclass[10pt]{article}
\usepackage[T1]{fontenc}
\usepackage[utf8]{inputenc}
\usepackage[french]{babel}
\usepackage{amsmath,amssymb,mathrsfs,mathpazo}
\usepackage{geometry}
\geometry{a4paper,margin=2cm}
\usepackage{graphicx,array,tabularx,tabularray,colortbl,multirow,multicol,booktabs,enumitem,pifont,ulem,fancybox,diagbox,cancel}
\usepackage{pstricks,pstricks-add,pst-plot,pst-tree,pst-node,pst-text,pst-func,pst-math,pst-bspline,pst-3dplot,pst-eucl,pst-solides3d}
\let\pscalculate\undefined
\usepackage{pst-bezier}
\usepackage{tikz,pgfplots,tkz-tab}
\usetikzlibrary{arrows,arrows.meta,patterns,calc,positioning,trees,decorations.markings,shapes.geometric,shapes.misc,shapes.arrows,intersections,decorations.pathmorphing,decorations.pathreplacing,backgrounds,shadows,babel}
\pgfplotsset{compat=1.18}
\usepgfplotslibrary{fillbetween}
\usepackage[np]{numprint}
\usepackage{listings,makecell,pas-tableur,twoopt}
\let\barre\undefined
\usepackage[pstricks]{tablvar}
\definecolor{darkgreen}{RGB}{0,100,0}
\providecommand{\red}{\color{red}}
\providecommand{\blue}{\color{blue}}
\providecommand{\green}{\color{green}}
\providecommand{\black}{\color{black}}
\providecommand{\white}{\color{white}}
\providecommand{\gray}{\color{gray}}
\providecommand{\cyan}{\color{cyan}}
\providecommand{\magenta}{\color{magenta}}
\providecommand{\yellow}{\color{yellow}}
\providecommand{\violet}{\color{violet}}
\providecommand{\orange}{\color{orange}}
\providecommand{\euro}{\texteuro}
\providecommand{\EUR}{\texteuro}
\providecommand{\ieme}{\textsuperscript{e}}
\providecommand{\degres}{\ensuremath{^\circ}}
\providecommand{\degree}{\ensuremath{^\circ}}
\providecommand{\vv}[1]{\overrightarrow{#1}}
\providecommand{\Coord}[1]{#1}
\providecommand{\decosix}{}
\providecommand{\starredbullet}{\ensuremath{\bullet}}
\pagestyle{empty}
\setlength{\parindent}{0pt}
\usepackage{adjustbox}\DeclareUnicodeCharacter{2236}{:}\DeclareUnicodeCharacter{2190}{\ensuremath{\leftarrow}}
\providecommand{\R}{\mathbb{R}}
\providecommand{\N}{\mathbb{N}}
\providecommand{\Z}{\mathbb{Z}}
\providecommand{\Q}{\mathbb{Q}}
\providecommand{\D}{\mathbb{D}}
\providecommand{\C}{\mathbb{C}}
\providecommand{\e}{{\mathrm{e}}}
\providecommand{\vect}[1]{\overrightarrow{#1}}
\providecommand{\vectt}[1]{\overrightarrow{\mathrm{#1}}}
\providecommand{\vvt}[1]{\overrightarrow{\mathrm{#1}}}
\providecommand{\np}[2][]{{#2~#1}}
\providecommand{\vv}[1]{{\overrightarrow{#1}}}
\renewcommand{\d}{{\mathrm{d}}}
\providecommand{\degres}{{^\circ}}
\providecommand{\degre}{{^\circ}}
\providecommand{\blue}{\color{blue}}
\providecommand{\red}{\color{red}}
\providecommand{\enskip}{\quad}
\providecommand{\llbracket}{\lbrack\!\lbrack}
\providecommand{\rrbracket}{\rbrack\!\rbrack}
\providecommand{\euro}{{\text{€}}}
\providecommand{\no}{{\text{n°}}}
\providecommand{\interval}[2]{{[#1;#2]}}
\providecommand{\intervFO}[2]{{[#1;#2[}}
\providecommand{\dotfill}{\ldots}

\providecommand{\ds}{\displaystyle}
\providecommand{\textding}[1]{#1}
\providecommand{\Oij}{\(\mathrm{O};\vec{\imath},\vec{\jmath}\)}
\providecommand{\Oijk}{\(\mathrm{O};\vec{\imath},\vec{\jmath},\vec{k}\)}
\providecommand{\Ouv}{\(\mathrm{O};\vec{u},\vec{v}\)}
\providecommand{\intervFF}[2]{\ensuremath{[#1;#2]}}
\renewcommand{\euro}{\ifmmode\text{\texteuro}\else\texteuro\fi}
\DeclareUnicodeCharacter{00B0}{\ifmmode^{\circ}\else\textdegree\fi}
\DeclareUnicodeCharacter{20AC}{\ifmmode\text{\texteuro}\else\texteuro\fi}
% Énoncés : Éducation nationale. Transcription et corrigés : APMEP, Denis Vergès et contributeurs. Découpage et classement : CoopMaths, Sébastien Lozano, Éric Elter et contributeurs. Adaptation de la présentation : Maths Explorer.
% CC BY-SA : https://coopmaths.fr/www/about/
% https://coopmaths.fr/alea/static/bac/2025/tex/bac_2025_09_sujet1_metropole_1.tex
\begin{document}
\noindent\textit{APMEP / CoopMaths. CC BY-SA. Présentation adaptée par Maths Explorer.}\par\medskip
\medskip

\textbf{Partie A}

\medskip

On considère l'équation différentielle

\[(E)\: \qquad y' + 0,4y = \e^{-0,4t}\]

où $y$ est une fonction de la variable réelle $t$.

On cherche l'ensemble des fonctions définies et dérivables sur $\R$ qui sont solutions
de cette équation.

\medskip

\begin{enumerate}
\item Soit $u$ la fonction définie sur $\R$ par : $u(t) = t\e^{-0,4t}$.

Vérifier que $u$ est solution de $(E)$.
\item Soit $f$ une fonction définie et dérivable sur $\R$.

On note $g$ la fonction définie sur $\R$ par : $g(t) = f(t) - u(t)$.

Soit $(H)$ l'équation différentielle $y' + 0,4y = 0$.
	\begin{enumerate}
		\item Démontrer que si la fonction $g$ est solution de l'équation différentielle $(H)$ alors la fonction $f$ est solution de l'équation différentielle $(E)$.
	\end{enumerate}		
		
On admettra que la réciproque est vraie.

	\begin{enumerate}[resume]
		\item Résoudre l'équation différentielle $(H)$.
		\item En déduire les solutions de $(E)$.
		\item Déterminer la solution $f$ de $(E)$ telle que $f(0) = 1$.
	\end{enumerate}
\end{enumerate}

\bigskip

\textbf{Partie B}

\medskip

On s'intéresse à la glycémie chez une personne venant de prendre un repas.

La glycémie en g.L$^{-1}$, en fonction du temps $t$, exprimé en heure, écoulé depuis la fin
du repas, est modélisée par la fonction $f$ définie sur [0~;~6] par :

\[f(t) = (t + 1)\e^{-0,4t}.\]

\medskip

\begin{enumerate}
\item
	\begin{enumerate}
		\item Montrer que, pour tout $t \in [0~;~6],\: f'(t) = (- 0,4t + 0,6)\e^{-0,4t}$.
		\item Étudier les variations de $f$ sur [0~;~6] puis dresser son tableau de variations sur cet intervalle.
	\end{enumerate}
\item Une personne est en hypoglycémie lorsque sa glycémie est inférieure à 0,7~g.L$^{-1}$.
	\begin{enumerate}
		\item Démontrer que sur l'intervalle [0~;~6] l'équation $f(t) = 0,7$ admet une unique
solution que l'on notera $\alpha$.
		\item Au bout de combien de temps après avoir pris son repas cette personne est-
elle en hypoglycémie ?

On exprimera ce temps à la minute près.
	\end{enumerate}
\item On souhaite déterminer la glycémie moyenne en g.L$^{-1}$ chez cette personne lors des six heures qui suivent le repas.
	\begin{enumerate}
		\item À l'aide d'une intégration par parties, montrer que:

\[\displaystyle\int_0^6 f(t)\: \text{d}t = - 23,75\e^{-2,4} + 8,75.\]

		\item Calculer la glycémie moyenne en g.L$^{-1}$ chez cette personne lors des six heures qui suivent le repas.
		\item En remarquant que la fonction $f$ est solution de l'équation différentielle $(E)$, expliquer comment on aurait pu obtenir ce résultat autrement.
	\end{enumerate}
\end{enumerate}

\bigskip

\end{document}
